% page 620 du fac-simile (image p-619 du scan)
% bloc 152.4 x 254.9 mm ; marge gauche 8.0 mm
\pgc{-12.700mm}{0.000mm}{
\l{\cel{2}612}
\l{}
\l{}
\l{venjar. (trans., su, ulu, pri, pro). Punisar ulu (qua}
\l{\cel{3}ofensis) por satisfacar sua rankoro, o la rankoro}
\l{\cel{3}di ulu. - EFIS.}
\l{}
\l{\sou{ventar}. (\sou{netrans.}) Dicesas kande, en ta o ca parto di}
\l{\cel{3}atmosfero, eventas diplaseso plu o min rapida di aero-}
\l{\cel{3}quanto. - DEFIS.}
\l{}
\l{\sou{ventilar}. (\sou{trans.}) Chanjar la aero di klozajo (+kluzajo),}
\l{\cel{3}aerizar per eventigar fluo di aero. - DEFIRS.}
\l{}
\l{\sou{ventilatoro}. (\sou{tekn.)}\cel{2}Aparato uzata por ventilar la cham-}
\l{\cel{3}brego di teatro, chambro di hospitalo, e c., o por ali-}
\l{\cel{3}mentar fornego ye aero. - EF.}
\l{}
\l{\sou{ventro}. (\sou{anat.)}\cel{2}Parto di la korpo, kava, qua kontenas la}
\l{\cel{3}stomako e la intestini. (sinonimo : abdomino). - EFIS.}
\l{}
\l{\sou{ventrikulo}. (\sou{anat.)}\cel{2}Kavajo qua sequas la parto supra di}
\l{\cel{3}la dextrajo e sinistrajo di la kordio ; la dextrala}
\l{\cel{3}sendas la sango veinala a la pulmoni - la sinistrala,}
\l{\cel{3}la sango arteriala, a la tota korpo. - DEFIS.}
\l{}
\l{\sou{ventuzo}. (\sou{medic.}) Mikra vazo klosho-forma, quan onu apli-}
\l{\cel{3}kas sur ta o ca parto di la korpo e de qua onu ekpulsas}
\l{\cel{3}la aero, por sublevar la pelo e produktar iriteso lokala}
\l{\cel{3}deturniva. - FIRS.}
\l{}
\l{\sou{\textquotedbl{}Venus\textquotedbl{}}. I. I. (\sou{astron.})\cel{2}Planeto di la sistemo Sunala, la}
\l{\cel{3}disto-duesma, kontante de Suno de qua lu distas per (cir-}
\l{\cel{3}kume) 108 milion kilometri. - II. (\sou{mitol}.) (\sou{che la Lati}-}
\l{\cel{3}ni). Deino di Amoro.- DEFIRS.}
\l{}
\l{\sou{venusmonto}. (\sou{anat.}) Saliajo avan la pubio,\cel{2}super la vulvo,}
\l{\cel{3}e qua pilozeskas lor pubereso. F.}
\l{}
\l{\sou{vera}. Konforma a to quo esas reala (kontraste kun : imituro,}
\l{\cel{3}kontrafakturo, finguro). EFIS.}
\l{}
\l{\sou{verando}. (\sou{arkit}.)\cel{2}Teraso tektoza qua formacas peristilo.-}
\l{\cel{3}DEFIR.}
\l{}
\l{\sou{verbo.}\cel{2}I. (\sou{logiko}). Termino di la propoziciono, qua unio-}
\l{\cel{3}nas la atributo a la subjekto. - II. (\sou{gram.}) Vorto-speco}
\l{\cel{3}qua expresas la existo, la ago. - II. (\sou{teol.}) Persono}
\l{\cel{3}duesma ek la Triuno kristana ; la filiulo di Deo, egar-}
\l{\cel{3}data kom la sajo eterna. - DEFIS.}
\l{}
\l{\sou{verbasko}.\cel{2}(\sou{bot.})\cel{2}Herboro moligiva, uzata por kataplasmi,}
\l{\cel{3}alta-statura, kun folii simpla, alterna, kovrita ye pili}
\l{\cel{3}kotonatra. - L.\cel{1}\sou{verbascum.}\cel{1}IL.}
\l{}
\l{\sou{verbeno}. (\sou{bot.)}\cel{1}Herboro di qua la stipito quadrata, la folii}
\l{\cel{3}inter-opozata e dentoza, la infloresenco spikoforma,}
\l{\cel{3}memorigas la labiacei - uzata por la faco di parfumi o}
\l{\cel{57}...}
}
